\chapter{Methodology} %This is an example title name

\begin{chapterabstract}

This is where you will insert your chapter abstract, which funcitons as a concise, self-contained summary of your chapter.

%Note that chapter abstracts are not required by KU. If your committee feels like you don't need an abstract for each chapter, simply don't include the \begin{chapterabstract} and \end{chapterabstract} commands in each chapter .tex file.

\end{chapterabstract}

%%%%%%%%%%%%%%%%%%%%%%%%%%%%%%%%%%%%%%%%%%%%%%%%%%%%%%%%%%%%%%%%%%%%%%%%%%%%%%%%%%%%%%%%%%
%%%%%%%%%%%%%%%%%%%%%%%%%%%%%%%%%%%%%%%%%%%%%%%%%%%%%%%%%%%%%%%%%%%%%%%%%%%%%%%%%%%%%%%%%%
% A section discussing the framework; i.e. the structure and
% flow of Dog treats
% A section discussing our application to Rust and Verus
% A section discussing how the dataset was found and constructed
% A Section the choice and use of LLM models


\section{DogTreats Framework}
% Summarize the purpose of dogtreats, the idea behind
Dogtreats is a method for effectively utilizing LLM's to generate the content of a proof, in order to automate part of the theorem proving process. However using LLM's alone to synthesize a full proof 
means that we could fail to reference objects in the local context that are either necessary or helpful for solving the proof. Furthermore in the case of failure diagnosing why the proof failed or
finding a partial solution from the provided proof can be difficult. Thus our method categorizes and ranks existing local information by relevance to provide to the LLM, and then attempts to iteratively
add and test single statements to the proof effectively constructing a partial solution at each step. The full process is given as follows: [How to transition into the full description?]

\subsection{Knowledge Graph Generation}
A knowledge graph for our purposes is a graph structure that relates an object $O_1$ to an object $O_2$ 

[Would a mathematical model of a knowledge graph be of benefit here]

\subsection{Error Tracking}

\subsection{Prompt Generation}

\subsection{Repair Attempt}

\section{Application to Verus}
Here is another example of a chapter subsection. Delete or duplicate as necessary.

\subsection{Selecting a Large Language Model}
Here is another example of a third-level subsection. Delete or duplicate as necessary.

\subsubsection{Dataset Construction}
Here is another example of a fourth-level subsection. Delete or duplicate as necessary.